\subsection{局部连通空间上局部常值层的态射}

设 B 为拓扑空间，\((\mathrm{E},p)\) 与 \((\mathrm{E}^{\prime},p^{\prime})\) 为 B-空间。用 \(\mathcal{M} = \underline{\mathcal{M}or_{\mathrm{B}}(\mathrm{E};\mathrm{E}^{\prime})}\) 表示从 \((\mathrm{E},p)\) 到 \((\mathrm{E}^{\prime},p^{\prime})\) 中的 B-态射在 B 上所成的层 (I, p. 45, example 4)。若 U 是 B 的开集而 b 是 U 的一点，我们用 \(\theta_{b,\mathrm{U}}\colon\mathcal{M}(\mathrm{U})\to\mathcal{C}(\mathrm{E}_{b};\mathrm{E}_{b}^{\prime})\) 表示通过过渡到 b 处的纤维而得到的典范映射。我们用 \(\theta_{b}\colon\mathcal{M}_{b}\to\mathcal{C}(\mathrm{E}_{b};\mathrm{E}_{b}^{\prime})\) 表示唯一的映射，使得对 B 的每个含 b 的开集 U，\(\theta_{b,\mathrm{U}}\) 都是 \(\theta_{b}\) 与典范映射 \(\mathcal{M}(\mathrm{U})\to\mathcal{M}_{b}\) 的复合 (E, III, p. 62)。

再设 \(\mathcal{I} = \underline{\mathcal{I}som}_{\mathrm{B}}(\mathrm{E};\mathrm{E}^{\prime})\) 为从 \((\mathrm{E},p)\) 到 \((\mathrm{E}^{\prime},p^{\prime})\) 中的 B-同构在 B 上所成的层。用 \(i\colon \mathcal{I}\to \mathcal{M}\) 表示典范态射。对每个 \(b\in \mathrm{B}\)，映射 \(\theta_{b}\circ i_{b}\) 诱导映射 \(\theta_b^\prime\) ，它把 \(\mathcal{I}_b\) 映到由 \(\mathrm{E}_b\) 到 \(\mathrm{E}_b^\prime\) 上的连续双射所成之集中。

命题 12. --- 设空间 B 是局部连通的且 B-空间 E 与 E' 是覆叠。于是层 \(\mathcal{M}or_{\mathrm{B}}(\mathrm{E};\mathrm{E}')\) 是局部常值的，且对每个 \(b\in\mathrm{B}\)，映射 \(\theta_{b}\) 是从其在 \(b\) 处的茎到从 \(\mathrm{E}_{b}\) 到 \(\mathrm{E}_{b}^{\prime}\) 中的映射所成之集上的双射。

类似地，层 \(\mathcal{I}\text{som}_{\mathrm{B}}(\mathrm{E};\mathrm{E}^{\prime})\) 是局部常值的，且对 B 的每一点 \(b\)，映射 \(\theta_{b}^{\prime}\) 是从其在 \(b\) 处的茎到从 \(\mathrm{E}_{b}\) 到 \(\mathrm{E}_{b}^{\prime}\) 上的双射所成之集上的双射。

我们将用 \(\mathcal{M} = \underline{\mathcal{M}or}_{\mathrm{B}}(\mathrm{E};\mathrm{E}^{\prime})\) 与 \(\mathcal{I} = \underline{\mathcal{Isom}}_{\mathrm{B}}(\mathrm{E};\mathrm{E}^{\prime})\) 表示。要证的断言在 B 上是局部性的；因此可设 \(\mathrm{E} = \mathrm{B}\times \mathrm{F}\) 与 \(\mathrm{E}' = \mathrm{B}\times \mathrm{F}'\) 是平凡覆叠，其中 F 与 \(\mathrm{F}'\) 是离散拓扑空间。

首先来证明，对 B 的每个连通开子集 U 与 U 的每一点 b，映射 \(\theta_{b,U}\) 是双射。设 \(f\colon U\times F\to U\times F'\) 为 U-空间的态射；于是 \(\theta_{b,U}(f)\) 是映射 \(y\mapsto\mathrm{pr}_{2}(f(b,y))\) （从 F 到 \(F'\) 中的）。对所有 \(y\in F\)，映射 \(x\mapsto\mathrm{pr}_{2}(f(x,y))\) 是从连通空间 U 到离散空间 \(F'\) 中的连续映射，故为常值的，等于 \(\mathrm{pr}_{2}(f(b,y))\)。于是有 \(f(x,y)=(x,\theta_{b,U}(f)(y))\)。由此得出 \(\theta_{b,U}\) 是双射；

其逆双射把每个映射 \(\varphi\colon\mathrm{F}\to\mathrm{F}^{\prime}\) 对应到从 \(\mathrm{U}\times\mathrm{F}\) 到 \(\mathrm{U}\times\mathrm{F}^{\prime}\) 中的、由 \((x,y)\mapsto(x,\varphi(y))\) 给出的 U-态射。

按假设，每点 \(b \in B\) 都有由连通开邻域组成的基本邻域系；于是映射 \(\theta_{b}\) 是双射。诸映射 \(\theta_{b,\mathrm{U}}^{-1}\) （其中 U 取 B 的非空连通开子集而 b 取 B 的任意点）定义了从以 \(F^{\prime F}\) 为茎型的常值层到层 \(\mathcal{M}\) 的一个预层态射（关于 B 的连通开子集组成的基）。由前所证，这个态射诱导茎上的双射；因而它是同构。

关于层 \(\mathcal{I}\) 的断言类似地证明。

推论 1. --- 设 B 为拓扑空间，A 为 B 的子空间。设空间 A 与 B 都是局部连通的。设 E 与 E' 为 B 的覆叠。则典范态射 \(\psi\colon\mathcal{M}or_{B}(E;E')_{A}\to\mathcal{M}or_{A}(E_{A};E'_{A})\) 与 \(\psi'\colon\mathcal{I}som_{B}(E;E')_{A}\to\mathcal{I}som_{A}(E_{A};E'_{A})\) (I, p. 45, example 4) 都是同构。

对每点 \(a \in \mathrm{A}\)，由前一个命题以及 I, p. 63 的例 2（应用于空间 \(\{a\}\)、A、B 及典范单射）得出，典范态射 \(\psi\) 通过过渡到茎诱导 \(\mathrm{E}_a^{\mathrm{E}_a'}\) 上的恒同映射。因而它是同构。\(\psi'\) 是同构这一事实类似地证明。

推论 2. --- 设 B 为拓扑空间，A 为 B 的子空间。设空间 A 与 B 都是局部连通的，且偶对 (B, A) 满足 I, p. 37 的性质 (PCV)。设 E 与 E' 为 B 的覆叠，用 p 与 p' 表示它们的投影。设 \(g: \overline{p}^{1}(A) \to (\overline{p}')^{1}(A)\) 为 A-态射（相应地，A-同构）。则存在 A 在 B 中的邻域 U 以及 U-态射（相应地，U-同构）\(f: \overline{p}^{1}(U) \to (\overline{p}')^{1}(U)\)，使得 \(f_{A} = g\)。

保留推论 1 的记号。由同一推论，A-态射 \(g \colon \mathrm{E}_{\mathrm{A}} \to \mathrm{E}_{\mathrm{A}}^{\prime}\) 与截面 \(s_0\) （平展空间 \(\mathrm{E}_{\mathcal{M}}\) （与 \(\mathcal{M}\) 相伴的）在 A 上的）相恒同。由对偶对 (B, A) 所作的假设以及 I, p. 39 的引理 3，存在 A 在 B 中的开邻域 U 以及连续截面 \(s\) （\(\mathrm{E}_{\mathcal{M}}\) 在 U 上的、延拓 \(s_0\) 的）。这个截面 \(s\) 与一个 U-态射 \(f \colon \mathrm{E}_{\mathrm{U}} \to \mathrm{E}_{\mathrm{U}}^{\prime}\) （延拓 \(g\) 的）相恒同。

g 是 A-同构的情形类似处理，只需把层 M 换成层 S。

